\section{Urysohn 度量化定理}\label{section:ch04-urysohn-metrization}

现在我们来证明本章的标志性定理，这个定理给出了拓扑空间为可度量化空间的条件。它的证明融汇了本书前面陈述的许多内容，不仅要用到第 2 章关于度量拓扑的结论，而且还要用到本章刚刚证明的与可数性公理及分离性公理有关的结论。证明的基本思路虽然简单但却很实用，读者将在后文中多次见到它的不同表达形式。

这里有两种证明方式，因为每一种方式在今后都有相应的推广，因此，我们把这两种方式都写出来。第一种方式将在第 5 章中证明关于完全正则空间的一个嵌入定理时得到推广，第二种方式则在第 6 章中证明 Nagata-Smirnov 度量化定理时得到推广。

\begin{theorem}[Urysohn 度量化定理]\label{theorem:ch04-urysohn-metrization}
每一个具有可数基的正则空间 $X$ 都是可度量化的。
\end{theorem}

\begin{proof}
我们通过证明 $X$ 能够被嵌入到一个可度量化的空间 $Y$ 中，也就是证明 $X$ 同胚于 $Y$ 的一个子空间，从而得到 $X$ 是可度量化的。两种证明的差别就在于对可度量化空间 $Y$ 的选择。第一种方式，将 $Y$ 取作具有积拓扑的空间 $\mathbb{R}^{\omega}$，我们已经证明了这个空间是可度量化的（见第 20 节）。第二种方式仍取 $Y$ 为 $\mathbb{R}^{\omega}$，但拓扑是由一致度量 $\overline{\rho}$（见第 20 节）所诱导的。在每一种情形下，我们的证明本质上都是要把 $X$ 嵌入到 $\mathbb{R}^{\omega}$ 的子空间 $[0,1]^{\omega}$ 中。

第一步。我们证明：存在由连续函数 $f_{n}\colon X\to[0,1]$ 构成的一个可数族，满足条件：对于 $X$ 中任意给定的一个点 $x_{0}$ 及 $x_{0}$ 的一个邻域 $U$，存在一个指标 $n$，使得 $f_{n}$ 在 $x_{0}$ 处取正值，而在 $U$ 的外部蜕化。

这是 Urysohn 引理的一个直接结果，即对于给定的 $x_{0}$ 和 $x_{0}$ 的一个邻域 $U$，存在这样一个连续函数。然而，倘若我们对于每一偶对 $(x_{0},U)$ 都选取一个相应的函数的话，那么所得到的函数族一般来说是不可数的。因此，我们的任务就是减小这个函数族的基数。以下是一种可行的方法：

设 $\{B_{n}\}$ 是 $X$ 的一个可数基。对于每一对使得 $\overline{B}_{n}\subset B_{m}$ 的指标 $n,m$，应用 Urysohn 引理，选取一个连续函数 $g_{n,m}\colon X\to[0,1]$ 使得 $g_{n,m}(\overline{B}_{n})=\{1\}$ 且 $g_{n,m}(X-B_{m})=\{0\}$。于是，函数族 $\{g_{n,m}\}$ 满足：对于给定的 $x_{0}$ 及 $x_{0}$ 的邻域 $U$，我们可以选取基中一个元素 $B_{m}$ 包含 $x_{0}$ 并且包含于 $U$。根据正则性，我们然后可以选取 $B_{n}$ 使得 $x_{0}\in B_{n}$ 和 $\overline{B}_{n}\subset B_{m}$。于是，对于指标的偶对 $n,m$，就给出了 $g_{n,m}$ 的定义，并且 $g_{n,m}$ 在 $x_{0}$ 取正值，在 $U$ 外蜕化。由于函数族 $\{g_{n,m}\}$ 的指标集是 $\mathbb{Z}_{+}\times\mathbb{Z}_{+}$ 的一个子集，所以它是一个可数族。因此它可以用正整数作为指标重新标记，于是我们便得到了所需的函数族 $\{f_{n}\}$。

第二步（第一种方式）。设 $\{f_{n}\}$ 为第一步中给出的函数族。赋予 $\mathbb{R}^{\omega}$ 积拓扑，并且给定一个映射 $F\colon X\to\mathbb{R}^{\omega}$，定义为：
\begin{align*}
F(x)=(f_{1}(x),f_{2}(x),\cdots).
\end{align*}
我们断言 $F$ 是一个嵌入。
\end{proof}

首先，由于 $\mathbb{R}^{\omega}$ 上具有积拓扑，并且每一个 $f_{n}$ 都是连续的，所以 $F$ 是连续的。其次，因为对于给定的 $x\neq y$，便存在一指标 $n$ 使得 $f_{n}(x)>0$ 和 $f_{n}(y)=0$，所以 $F(x)\neq F(y)$。因此 $F$ 是一个单射。

最后，我们还须证明 $F$ 是从 $X$ 到像集的一个同胚，也就是到 $\mathbb{R}^{\omega}$ 的子空间 $Z=F(X)$ 的一个同胚。我们知道：$F$ 定义了一个从 $X$ 到 $Z$ 的连续的一一映射，因而仅需证明对于 $X$ 的每一开集 $U$，$F(U)$ 是 $Z$ 中的一个开集。设 $z_{0}$ 为 $F(U)$ 中一点，我们只需证明存在 $Z$ 中的一个开集 $W$ 使得
\begin{align*}
z_{0}\in W\subset F(U).
\end{align*}

设 $x_{0}$ 为 $U$ 的一个点，使得 $F(x_{0})=z_{0}$。选取一个指标 $N$ 使得 $f_{N}(x_{0})>0$ 并且 $f_{N}(X-U)=\{0\}$。取 $\mathbb{R}$ 中的开射线 $(0,+\infty)$，令 $V$ 表示 $\mathbb{R}^{\omega}$ 中的开集
\begin{align*}
V=\pi_{N}^{-1}((0,+\infty)).
\end{align*}
设 $W=V\cap Z$。于是根据子空间拓扑的定义可见 $W$ 是 $Z$ 的一个开集。参见\cref{figure:ch04-embedding-rw}。我们断言：$z_{0}\in W\subset F(U)$。首先，$z_{0}\in W$，因为
\begin{align*}
\pi_{N}(z_{0})=\pi_{N}(F(x_{0}))=f_{N}(x_{0})>0.
\end{align*}

\par
\begin{figure}[htbp]
\centering
\includegraphics[width=0.6\textwidth]{ch04-fig34-1.png}
\caption{嵌入到 $\mathbb{R}^{\omega}$}\label{figure:ch04-embedding-rw}
\end{figure}

其次，$W\subset F(U)$。这是由于：若 $z\in W$，则对于某一个 $x\in X$ 有 $z=F(x)$，并且 $\pi_{N}(z)\in(0,+\infty)$。因为 $\pi_{N}(z)=\pi_{N}(F(x))=f_{N}(x)$ 以及 $f_{N}$ 在 $U$ 外蜕化，所以点 $x$ 必定属于 $U$。于是 $z=F(x)$ 在 $F(U)$ 中。这便是我们所要证明的。

从而 $F$ 便是一个 $X$ 到 $\mathbb{R}^{\omega}$ 的嵌入。

第三步（第二种方式）。在以下证明中，我们是将 $X$ 嵌入到度量空间 $(\mathbb{R}^{\omega},\overline{\rho})$ 中。事实上，我们是将 $X$ 嵌入到度量空间 $([0,1]^{\omega},\overline{\rho})$ 中，其中度量 $\overline{\rho}$ 就等于度量
\begin{align*}
\rho(\boldsymbol{x},\boldsymbol{y})=\sup\{\,|x_{i}-y_{i}|\,\}.
\end{align*}
我们应用第一步中构造的函数 $f_{n}\colon X\to[0,1]$ 的可数族。但这次我们附加一个条件：对于所有的 $x$，$f_{n}(x)\leqslant 1/n$。（这个条件是容易满足的，我们仅需将每一个函数 $f_{n}$ 除以 $n$ 便可以了。）

像前面一样，定义函数 $F\colon X\to[0,1]^{\omega}$ 使其满足
\begin{align*}
F(x)=(f_{1}(x),f_{2}(x),\cdots).
\end{align*}
我们断言：$F$ 是到度量空间 $([0,1]^{\omega},\rho)$ 中的一个嵌入。由第二步的证明中可见 $F$ 是一个单射。此外，我们知道：若在 $[0,1]^{\omega}$ 上取积拓扑，那么映射 $F$ 将 $X$ 的开集映为子空间 $Z=F(X)$ 的开集。而当我们考虑由度量 $\rho$ 诱导的 $[0,1]^{\omega}$ 上更细的拓扑时，上述性质依然成立。

余下有待证明的就是 $F$ 的连续性。它已不再是各个分量函数连续性的直接推论，因为现在使用的已经不再是 $\mathbb{R}^{\omega}$ 上的积拓扑了。这就是要附加 $f_{n}(x)\leqslant 1/n$ 这个条件的原因。

设 $x_{0}$ 是 $X$ 的一个点，$\varepsilon>0$。为了证明连续性，我们需要找到 $x_{0}$ 的一个邻域 $U$ 使得对于所有 $x\in U$，有
\begin{align*}
x\in U\Longrightarrow\rho(F(x),F(x_{0}))<\varepsilon.
\end{align*}
首先，选取充分大的 $N$，使得 $1/N\leqslant\varepsilon/2$。那么对于每一个 $n=1,\cdots,N$ 应用 $f_{n}$ 的连续性，选取 $x_{0}$ 的一个邻域 $U_{n}$ 使得对于 $x\in U_{n}$ 有
\begin{align*}
|f_{n}(x)-f_{n}(x_{0})|\leqslant\varepsilon/2.
\end{align*}
令 $U=U_{1}\cap\cdots\cap U_{N}$。我们证明 $U$ 就是所求的 $x_{0}$ 的邻域。设 $x\in U$。若 $n\leqslant N$，则由 $U$ 的选取可见
\begin{align*}
|f_{n}(x)-f_{n}(x_{0})|\leqslant\varepsilon/2.
\end{align*}
若 $n>N$，则
\begin{align*}
|f_{n}(x)-f_{n}(x_{0})|<1/N\leqslant\varepsilon/2.
\end{align*}
因为 $f_{n}$ 将 $X$ 映入 $[0,1/n]$。所以，对于每一个 $x\in U$ 有
\begin{align*}
\rho(F(x),F(x_{0}))\leqslant\varepsilon/2<\varepsilon.
\end{align*}
定理证毕。

在上述证明的第二步中，实际上我们证明了一个比所陈述的结果更强的结论。为以后使用方便起见，我们在此将其表述如下：

\begin{theorem}[嵌入定理]\label{theorem:ch04-embedding}
设 $X$ 是一个空间，其中的每一个单点集是闭集。假定 $\{f_{\alpha}\}_{\alpha\in J}$ 是连续函数的一个加标族，其中 $f_{\alpha}\colon X\to\mathbb{R}$ 满足：对于 $X$ 中每一点 $x_{0}$ 以及 $x_{0}$ 的每一邻域 $U$，存在一个指标 $\alpha$ 使得 $f_{\alpha}$ 在 $x_{0}$ 取正值并且在 $U$ 之外取 $0$。那么由
\begin{align*}
F(x)=(f_{\alpha}(x))_{\alpha\in J}
\end{align*}
所定义的函数 $F\colon X\to\mathbb{R}^{J}$ 是一个 $X$ 到 $\mathbb{R}^{J}$ 的嵌入。若对于每一个 $\alpha$ 有 $f_{\alpha}$ 将 $X$ 映入 $[0,1]$ 中，则 $F$ 将 $X$ 嵌入 $[0,1]^{J}$。
\end{theorem}

其证明几乎完全是上述定理证明中第二步的翻版，只须用 $\alpha$ 代替 $n$，用 $\mathbb{R}^{J}$ 代替 $\mathbb{R}^{\omega}$。其中 $X$ 中的单点集是闭集的条件是为了保证当 $x\neq y$ 时存在一个指标 $\alpha$ 使得 $f_{\alpha}(x)\neq f_{\alpha}(y)$。

满足这一定理中条件的连续函数族称之为分离 $X$ 中的点与闭集的函数族。我们已清楚地看到：对于每一个满足单点集是闭集的空间，这样的函数族的存在等价于 $X$ 是完全正则的。因而，我们有以下直接推论：

\begin{theorem}\label{theorem:ch04-completely-regular-cube}
一个空间 $X$ 是完全正则空间当且仅当对于某一个 $J$，$X$ 同胚于 $[0,1]^{J}$ 的一个子空间。
\end{theorem}

\subsection*{习题}

\begin{enumerate}
\item 举例说明具有可数基的 Hausdorff 空间未必是一个可度量化空间。
\item 举例说明一个完全正规并且满足第一可数性公理、Lindelöf 条件及具有可数稠密子集的空间未必是一个可度量化空间。
\item 设 $X$ 是一个紧致的 Hausdorff 空间。证明 $X$ 是可度量化的当且仅当 $X$ 有可数基。
\item 设 $X$ 是一个局部紧致的 Hausdorff 空间。$X$ 有可数基是否意味着 $X$ 是可度量化的？$X$ 是一个度量空间是否意味着 $X$ 有可数基？
\item 设 $X$ 是一个局部紧致的 Hausdorff 空间。设 $Y$ 为 $X$ 的单点紧致化。$X$ 有可数基是否蕴涵着 $Y$ 是可度量化的？$Y$ 可度量化是否意味着 $X$ 有可数基？
\item 验证\cref{theorem:ch04-embedding} 证明的细节。
\item 一个空间称为\textbf{局部可度量化的}（locally metrizable），如果对于 $X$ 的每一个点 $x$，存在一个邻域关于子空间拓扑是可度量化的。证明：一个紧致的 Hausdorff 空间，如果它是局部可度量化的，那么它是可度量化的。[提示：证明 $X$ 是有限个具有可数基的开子空间之并。]
\item 证明：一个局部可度量化空间，如果它是正则的 Lindelöf 空间，那么它是可度量化的。[提示：Lindelöf 空间的闭子空间是 Lindelöf 空间。]正则性是关键的。你的证明中何处用到了它？
\item 设紧致的 Hausdorff 空间 $X$ 是两个闭子空间 $X_{1}$ 和 $X_{2}$ 之并。若 $X_{1}$ 和 $X_{2}$ 是可度量化的，证明 $X$ 也是可度量化的。[提示：构造 $X$ 的开集所组成的一个可数族 $\mathcal{A}$，使得对于 $i=1,2$，$\mathcal{A}$ 的元素与 $X_{i}$ 的交构成了 $X_{i}$ 的一个基。假设 $X_{1}-X_{2}$ 和 $X_{2}-X_{1}$ 属于 $\mathcal{A}$。取 $\mathcal{B}$ 为 $\mathcal{A}$ 中所有元素的有限交所构成的族。]
\end{enumerate}